\documentclass[11pt,a4paper,oneside]{book}

% ---------- Encoding & fonts ----------
\usepackage[T1]{fontenc}
\usepackage[utf8]{inputenc}
\usepackage{lmodern}
\usepackage[english]{babel}
\usepackage{microtype}

% ---------- Page layout ----------
\usepackage[margin=2.5cm]{geometry}
\usepackage{fancyhdr}
\pagestyle{fancy}
\fancyhf{}
\fancyhead[L]{\nouppercase{\leftmark}}
\fancyhead[R]{\thepage}
\renewcommand{\headrulewidth}{0.4pt}

% ---------- Math ----------
\usepackage{amsmath,amssymb,amsthm}
\usepackage{bm}

% ---------- Graphics, tables, code ----------
\usepackage{graphicx}
\graphicspath{{figures/}}
\usepackage{booktabs}
\usepackage{listings}
\usepackage{xcolor}

\lstdefinestyle{pystyle}{
  language=Python,
  basicstyle=\ttfamily\small,
  keywordstyle=\color{blue!70!black},
  commentstyle=\color{gray},
  stringstyle=\color{green!50!black},
  numbers=left,
  numberstyle=\tiny\color{gray},
  frame=single,
  breaklines=true,
  showstringspaces=false,
}
\lstset{style=pystyle}

% ---------- References ----------
\usepackage{hyperref}
\hypersetup{
  colorlinks=true,
  linkcolor=blue!50!black,
  citecolor=blue!50!black,
  urlcolor=blue!50!black,
}
\usepackage{cleveref}

% ---------- Theorem environments ----------
\theoremstyle{definition}
\newtheorem{definition}{Definition}[chapter]
\newtheorem{example}[definition]{Example}
\newtheorem{exercise}[definition]{Exercise}

\theoremstyle{plain}
\newtheorem{theorem}[definition]{Theorem}
\newtheorem{proposition}[definition]{Proposition}

\theoremstyle{remark}
\newtheorem*{remark}{Remark}

% ---------- Custom macros ----------
\newcommand{\R}{\mathbb{R}}
\newcommand{\vect}[1]{\bm{#1}}
\newcommand{\deriv}[2]{\frac{\partial #1}{\partial #2}}
\DeclareMathOperator*{\argmin}{arg\,min}

% ---------- Title ----------
\title{\Huge Machine Learning for Physicists}
\author{Caio Lagana}
\date{\today}

\begin{document}

\frontmatter
\maketitle
\tableofcontents

\chapter{Preface}

So, you're a physicist looking to dive into the world of machine learning. You have a deep understanding of calculus, are familiar with experimental data-fitting methods, and have mastered concepts like entropy and non-linearity---yet you have only a vague idea of how a machine learning algorithm is actually written and how it works? Then this book is for you. And the good news is: your learning curve will grow fast. I've been in your shoes. I used to view machine learning as a distant horizon, but I soon realized that the knowledge I'd gained in physics made it incredibly easy to quickly master the concepts and engineering behind the field. So, come along and be amazed.

\mainmatter

\chapter{Single-neuron ML}
\label{chap:single-neuron}

You have been using a single-neuron network since the first time you fitted experimental data, but perhaps no one ever told you that. Let's recall the basics. You have a set of experimental data points
\begin{equation}
    \label{eq:xy-dataset}
    (x_i,y_i)
\end{equation}
through which to fit a linear function
\begin{equation}
    \label{eq:linear-activation}
    f(x)=ax+b.
\end{equation}

The Mean Squared Error, defined as
\begin{equation}
    \begin{split}
        \label{eq:mse-loss}
        \text{MSE}(a,b) &= \frac{1}{N}\sum_{i=1}^N[y_i - f(x_i)]^2 \\
                        &= \frac{1}{N}\sum_{i=1}^N[y_i - (ax_i + b)]^2
    \end{split}
\end{equation}
tells us how distant the curve $f(x)$ is from the datapoints. Fitting $f(x)$ to the datapoints $(x_i,y_i)$ means minimizing $\text{MSE}(a,b)$, that is, finding $a$ and $b$ such that
\begin{equation}
    \label{eq:single-neuron-backpropagation}
    \frac{\partial}{\partial(a,b)}\text{MSE}(a,b) = 0.
\end{equation}

No surprises so far, right? Well, what if I tell you the above is the simplest possible neural network? Indeed, it is a single-neuron network, where \eqref{eq:linear-activation} is the activation function, \eqref{eq:mse-loss} is the loss function and \eqref{eq:single-neuron-backpropagation} is a very rudimentary backpropagation method.

Solving \eqref{eq:single-neuron-backpropagation} for $a$ and $b$, which, in this simple case is possible analytically, fixes the parameters to $a_0$ and $b_0$
\begin{equation}
    \begin{split}
        \label{eq:a0b0}
        a_0 &= \frac{\sum_i(x_i-\bar x)(y_i-\bar y)}{\sum_i(x_i-\bar x)^2} \\
        b_0 &= \bar y-a\bar x
    \end{split}
\end{equation}
and we end up with a model $f_0(x) = a_0x+b_0$ (in \eqref{eq:a0b0}, $\bar x$ and $\bar y$ are the mean values of the $x_i$ and $y_i$ datapoints). This model is a single neuron, no hidden layers: it receives one input parameter and outputs a number. $f_0$ has learned from data \eqref{eq:xy-dataset} through \eqref{eq:single-neuron-backpropagation} and can now predict on any given $x$. Pretty cool, huh? The concepts were sitting there all the time. Take a moment to re-read the definitions and let's jump to the second simplest model: a two-neuron network.


\chapter{Two-neuron ML}


Let us now concatenate two neurons.


\backmatter

\bibliographystyle{plain}
\bibliography{references}

\end{document}
